% Zdroj: PDF priklady-leto-2022-one.pdf (sešit S1, PGVVH+UI), fyzická str. 3. Značka: specializační otázka UI. Zařazeno: S3. Číselné výsledky ověřeny.
\section{Shlukování (K-means)}
\szzotazka{léto 2022}{12}{Shlukování (K-means)}

\noindent\textbf{Zadání.}\par
Ručně provedeme první iteraci algoritmu K-means pro $K = 2$ na množině šesti příkladů s~příznaky
$A_1, A_2$:

\begin{center}
\begin{tabular}{ccc}
\toprule
příklad & $A_1$ & $A_2$ \\
\midrule
1 & 2 & 4 \\
2 & 1 & 3 \\
3 & 0 & 4 \\
4 & 5 & 2 \\
5 & 6 & 2 \\
6 & 4 & 0 \\
\bottomrule
\end{tabular}
\end{center}

Náhodně jsme příklady přiřadili do dvou shluků: $C_1 = \{2,3,4\}$, $C_2 = \{1,5,6\}$.
\begin{enumerate}
  \item Spočítejte pro každý shluk jeho centroid.
  \item Přiřaďte každý příklad k~nejbližšímu centroidu (euklidovská vzdálenost) a~uveďte prvky
    nově vzniklých shluků.
\end{enumerate}

\medskip
\noindent\textbf{Řešení.} \textit{(V~PDF bez náčrtu řešení; vlastní vzorové řešení, výpočet ověřen numericky.)}\par
\begin{enumerate}
  \item Centroid je průměr bodů shluku:
    \[
      c_1 = \Big(\tfrac{1+0+5}{3},\,\tfrac{3+4+2}{3}\Big) = (2,3),\qquad
      c_2 = \Big(\tfrac{2+6+4}{3},\,\tfrac{4+2+0}{3}\Big) = (4,2).
    \]
  \item Pro každý bod porovnáme čtverce vzdáleností k~$c_1=(2,3)$ a~$c_2=(4,2)$:
    \begin{center}
    \begin{tabular}{ccccc}
    \toprule
    bod & $(A_1,A_2)$ & $\|\cdot - c_1\|^2$ & $\|\cdot - c_2\|^2$ & shluk \\
    \midrule
    1 & $(2,4)$ & $1$  & $8$  & $C_1$ \\
    2 & $(1,3)$ & $1$  & $10$ & $C_1$ \\
    3 & $(0,4)$ & $5$  & $20$ & $C_1$ \\
    4 & $(5,2)$ & $10$ & $1$  & $C_2$ \\
    5 & $(6,2)$ & $17$ & $4$  & $C_2$ \\
    6 & $(4,0)$ & $13$ & $4$  & $C_2$ \\
    \bottomrule
    \end{tabular}
    \end{center}
    Nově vzniklé shluky: $C_1 = \{1,2,3\}$, $C_2 = \{4,5,6\}$.
\end{enumerate}

